\section{Formal Model}
\label{sec:model-ra}

The model abstracts away the OS executor, the canonical JSON encoder,
and the receipt ledger; those lower layers are audited assumptions in
the proof manifest. The institutional core that remains is the
executive-action type, its TTL invariant, the rollback receipt, the
TTL-bounded amendment chain, and the composition theorem connecting
the response side to the parent paper's amendment side.

\paragraph{Executive action.}
An executive action is a structure that cannot be constructed without
a positive TTL witness:
\[
  \mathsf{ExecutiveAction} =
  (\mathit{rid},\,\mathit{t_{\text{issued}}},\,\mathit{ttl},\,
   \mathit{ttlPos},\,\mathit{kind},\,\mathit{rb})
\]
where $\mathit{rid}$ is the receipt identifier, $\mathit{t_{\text{issued}}}$
is the issued-at instant, $\mathit{ttl}$ is the duration in counter
units, $\mathit{ttlPos}$ is the field-level proof $0 < \mathit{ttl}$,
$\mathit{kind}$ is the action class tag, and $\mathit{rb}$ is an
optional rollback receipt. The expiry $\mathit{t_{\text{expires}}} =
\mathit{t_{\text{issued}}} + \mathit{ttl}$ is derived structurally.

The action is closed at instant $t$ iff a rollback receipt exists with
$\mathit{rolledBackAt} \le t$ or the TTL window has elapsed by $t$.
This is recorded as the definitional bridge
\thm{rollback_closes_or_ttl_window_active}, which discharges to
\emph{rfl} after case analysis on the rollback option. The bridge is
not the headline theorem; it documents the relationship between the
type-level fields and the propositional closure claim, in the same
role \thm{amendment_admissible_iff_backward_refinement} plays for the
amendment side of the parent paper.

\paragraph{TTL-bounded amendment.}
A TTL-bounded amendment $A$ is a constitutional delta $\delta$ paired
with an issued-at instant $\mathit{t_{\text{issued}}}$ and a positive
TTL duration $\mathit{ttl}$. The active constitution at instant $t$ is
\[
  \mathsf{activeAt}(A, t) =
  \begin{cases}
    \delta.\mathit{old} & \text{if } \mathit{t_{\text{issued}}} + \mathit{ttl} \le t, \\
    \delta.\mathit{new} & \text{otherwise.}
  \end{cases}
\]
The case split is a substrate-level abstraction of the TTL scheduler's
behavior. The scheduler advances the counter and reverts the amendment
at $\mathit{t_{\text{expires}}}$; the model represents this as a
conditional on $t$.

\paragraph{The composition theorem.}
The headline composition theorem
\thm{ttl_bounded_amendment_chain_preserves_baseline} states that a
chain of TTL-bounded amendments, each carrying a backward-refinement
witness against a fixed baseline constitution, preserves the baseline
at every instant:
\[
  \forall t,\, \forall A \in \mathit{chain},\
    \mathsf{BackwardRefines}\bigl(\mathsf{activeAt}(A, t),\, \mathit{baseline}\bigr).
\]
The proof shape is structural rather than definitional. For each
amendment $A$ in the chain, the conclusion case-splits on whether $t$
falls inside or outside $A$'s window. Outside the window,
$\mathsf{activeAt}(A, t) = \delta.\mathit{old} = \mathit{baseline}$ and
the conclusion holds by reflexivity of $\mathsf{BackwardRefines}$ on a
fixed constitution. Inside the window,
$\mathsf{activeAt}(A, t) = \delta.\mathit{new}$ and the conclusion
holds by the per-step refinement witness on $A$.

The shape is the response-side dual of the parent paper's
\thm{essential_preserved_chain}, which composes per-step essential-
admission witnesses over an amendment sequence. The dual relation is
that the parent theorem ranges over receipts and a chain of
constitutions; this theorem ranges over instants and a chain of
TTL-bounded amendments. Both theorems are structurally inductive over
the chain rather than definitional on a single step. The discharge in
each case threads the per-step witness through the case analysis on
the inductive variable.

\paragraph{Rollback admissibility.}
The supporting reduction \thm{rollback_admission_composes_with_refinement}
states that a rollback receipt admissible under the post-amendment
constitution remains admissible under the pre-amendment constitution
given the refinement witness:
\[
  \mathsf{admits}(c_{\text{new}}, \mathit{rb}) = \text{true}
    \;\wedge\; \mathsf{BackwardRefines}(c_{\text{new}}, c_{\text{old}})
    \;\Rightarrow\; \mathsf{admits}(c_{\text{old}}, \mathit{rb}) = \text{true}.
\]
The proof is a direct application of the refinement hypothesis to the
rollback receipt's admission. The theorem's role in the paper is the
load-bearing reduction the headline composition leans on: it shows
that the rollback receipt produced inside an amendment's TTL window
does not become unadmittable when the window expires and the baseline
returns. In the substrate, this is the property the auto-reversion
relies on: the rollback receipt's canonical body is the same under
both constitutions, and the predicates the baseline applies to it are
a subset of the predicates the amended constitution applies.

\paragraph{Bilateral destructive admission.}
The bilateral reduction \thm{destructive_action_requires_bilateral_admission}
specializes the parent paper's
\thm{treaty_admission_iff_predicate_intersection} to the destructive
subclass of the action taxonomy. For a bilateral envelope binding a
device-polity receipt id and an operator-polity receipt id under a
destructive action class, admission requires that each polity's
predicates accept its own receipt. The reduction is structural rather
than novel: the parent theorem already establishes the bilateral
intersection; this theorem renames it for the destructive class. The
contribution at this point is the typed envelope, not a new theorem.

\paragraph{Threat model.}
The adversary can author and sign requests, attempt to skip the
rollback receipt, forge or delete the rollback receipt in the local
ledger, backdate or freeze the substrate counter, omit the
acknowledgement receipt, present a rollback receipt for an execution
it did not perform, or claim that an inverse executor completed when
it did not. The adversary cannot break Ed25519, SHA-256 collision
resistance, or canonical JSON correctness. The substrate counter
itself is an assumption: a compromised counter falsifies the TTL claim
in the same way a compromised signing key falsifies admission. The
counter is named as a foundational assumption in
Section~\ref{sec:limitations-ra}.

The model discharges two obligations. First, the typed executive-
action constructor refuses to build a witness without a positive TTL
and a class tag, so the type system rules out unbounded acts on the
construction side. Second, the composition theorem pins the chain's
trajectory to the baseline at every instant, so an amendment ratchet
that aggregates predicate drift through a sequence of TTL-bounded
amendments cannot accumulate without violating a per-step refinement
witness. The third obligation, namely that the rollback receipt
faithfully witnesses the inverse executor's completion, is a runtime
obligation discharged by the executor's signed completion against the
forward effect's content hash and is not a theorem in the bounded
model.
